%%
%% This is file `sample-sigconf-biblatex.tex',
%% generated with the docstrip utilŹty.
%%
%% The original source files were:
%%
%% samples.dtx  (with options: `all,proceedings,sigconf-biblatex')
%%
%% IMPORTANT NOTICE:
%%
%% For the copyright see the source file.
%%
%% Any modified versions of this file must be renamed
%% with new filenames distinct from sample-sigconf-biblatex.tex.
%%
%% For distribution of the original source see the terms
%% for copying and modification in the file samples.dtx.
%%
%% This generated file may be distributed as long as the
%% original source files, as listed above, are part of the
%% same distribution. (The sources need not necessarily be
%% in the same archive or directory.)
%%
%%
%% Commands for TeXCount
%TC:macro \cite [option:text,text]
%TC:macro \citep [option:text,text]
%TC:macro \citet [option:text,text]
%TC:envir table 0 1
%TC:envir table* 0 1
%TC:envir tabular [ignore] word
%TC:envir displaymath 0 word
%TC:envir math 0 word
%TC:envir comment 0 0
%%
%% The first command in your LaTeX source must be the \documentclass
%% command.
%%
%% For submission and review of your manuscript please change the
%% command to \documentclass[manuscript, screen, review]{acmart}.
%%
%% When submitting camera ready or to TAPS, please change the command
%% to \documentclass[sigconf]{acmart} or whichever template is required
%% for your publication.
%%
%%
% \documentclass[sigconf,natbib=false]{acmart}
\documentclass[sigconf,natbib=false]{acmart}
%%
%% \BibTeX command to typeset BibTeX logo in the docs
\AtBeginDocument{%
  \providecommand\BibTeX{{%
    Bib\TeX}}}

%%
%% For managing citations, it is recommended to use bibliography
%% files in BibTeX format.
%%
%% You can then either use BibTeX with the ACM-Reference-Format style,
%% or BibLaTeX with the acmnumeric or acmauthoryear sytles, that include
%% support for advanced citation of software artefact from the
%% biblatex-software package, also separately available on CTAN.
%%
%% Look at the sample-*-biblatex.tex files for templates showcasing
%% the biblatex styles.
%%

%%
%% The majority of ACM publications use numbered citations and
%% references, obtained by selecting the acmnumeric BibLaTeX style.
%% The acmauthoryear BibLaTeX style switches to the "author year" style.
%%
%% If you are preparing content for an event
%% sponsored by ACM SIGGRAPH, you must use the acmauthoryear style of
%% citations and references.
%%
%% Bibliography style
% \RequirePackage[
%   datamodel=acmdatamodel,
%   style=acmnumeric,
%   ]{biblatex}

% \RequirePackage[
% backend=biber,
% style=alphabetic,
% sorting=nyt
% ]{biblatex}

\RequirePackage[
  datamodel=acmdatamodel,
  style=acmnumeric,
  sorting=nyt
  ]{biblatex}

%% Declare bibliography sources (one \addbibresource command per source)
% \addbibresource{software.bib
% \addbibresource{sample-base.bib}
\addbibresource{papers.bib}

%% TODO support
\usepackage{xcolor}
\newcommand{\todo}[1]{\textcolor{red}{\textbf{[TODO: #1]}}}

%%
%% end of the preamble, start of the body of the document source.
\begin{document}

%%
%% The "title" command has an optional parameter,
%% allowing the author to define a "short title" to be used in page headers.
\title{ChronoQA: Benchmarking Temporal Grounding in Long-Form Video via Timestamp-Anchored Multiple-Choice QA}

%%
%% The "author" command and its associated commands are used to define
%% the authors and their affiliations.
%% Of note is the shared affiliation of the first two authors, and the
%% "authornote" and "authornotemark" commands
%% used to denote shared contribution to the research.
\author{Oluwatumininu Oguntola}
\email{oguntola@unc.edu}
\orcid{0009-0006-1999-1838}
\affiliation{%
  \institution{University of North Carolina at Chapel Hill}
  \city{Chapel Hill}
  \state{North Carolina}
  \country{USA}
}

\author{Pranav Wagh}
\email{pawagh@unc.edu}
\affiliation{%
  \institution{University of North Carolina at Chapel Hill}
  \city{Chapel Hill}
  \state{North Carolina}
  \country{USA}
}

\author{Gedas Bertasius}
\email{gedas@cs.unc.edu}
\affiliation{%
  \institution{University of North Carolina at Chapel Hill}
  \city{Chapel Hill}
  \state{North Carolina}
  \country{USA}
}

%%
%% By default, the full list of authors will be used in the page
%% headers. Often, this list is too long, and will overlap
%% other information printed in the page headers. This command allows
%% the author to define a more concise list
%% of authors' names for this purpose.
\renewcommand{\shortauthors}{Oguntola et al.}

%%
%% The abstract is a short summary of the work to be presented in the article.
%% Needs headline result once experiments run.
\begin{abstract}
The expanding use of large language models in intelligent workflows has triggered rapid advancement in Multimodal Large Language models (MLLM). Despite their success with short-form video, existing models still fall short when tasked with long-duration sequences.
Temporal grounding, which involves localizing when an event occurs within a video, is a persistent weakness of video-language models (VLMs), in part because existing training data provides answer supervision but lacks temporal evidence supervision. 
Without it, artificial intelligence cannot reliably assist with tasks such as medical surgery review, security monitoring, or analyzing long-form video data. Yet, existing benchmarks evaluate answer correctness without verifying whether a model identifies the correct supporting temporal evidence.
% Evaluating [N] VLMs reveals a consistent gap between QA accuracy and Mean Temporal IoU (mIoU), indicating that model frequently select correct answers without identifying the supporting temporal evidence.
We introduce \textbf{ChronoQA}, a multiple-choice QA benchmark that pairs every correct answer with a timestamp interval, providing an explicit grounding target that other video QA datasets omit, and enabling evaluation of answer accuracy and temporal localization quality.
We construct ChronoQA from [X] long-form YouTube videos (30-90 minutes each), yielding [Y] timestamp-anchored multiple choice QA pairs across five task categories: Action Understanding, Needle-in-the-Haystack, Ordering, Causal Understanding, and General. Through multimodel blind filtering, we constrain text-only solvability to [Z]\%, near the 20\% random baseline, ensuring visual perception rather than linguistic bias and visual information leakage. 
% Evaluating [N] VLMs reveals a consistent gap between QA accuracy and Mean Temporal IoU (mIoU), indicating that models frequently select correct answers without identifying the temporal evidence that supports them. ChronoQA provides a concrete diagnostic for this gap that existing benchmarks cannot measure.
% Models trained on ChronoQA can be evaluated on both answer accuracy and Mean Temporal IoU (mIoU), enabling direct measurement of whether training on grounded supervision improves temporal localization alongside answer quality.
\end{abstract}

%%
%% Keywords. The author(s) should pick words that accurately describe
%% the work being presented. Separate the keywords with commas.
% \keywords{Do, Not, Use, This, Code, Put, the, Correct, Terms, for,
%   Your, Paper}

% \received{16 March 2026}
% \received[revised]{12 March 2009}
% \received[accepted]{5 June 2009}

%%
%% This command processes the author and affiliation and title
%% information and builds the first part of the formatted document.
\maketitle

\section{Introduction}
% CITE EVERYTHING HERE, EVERY BENCHMARK, MODEL, ETC.
While conventional video models are proficient in determining the occurrence of an event, they frequently struggle with pinpointing the timing and duration of these events. A student reviewing a recorded lecture must locate the exact moment a concept is introduced, understand how it is developed over the rest of the video, and identify the follow-up examples that make it concrete enough to apply.

% A legal auditor reviewing a lengthy recording of a deposition must identify the exact timestamp when a key witness provides critical testimony, understand the duration of that testimony, and know the timestamps of the questions posed immediately before and after.

% In a clinical review setting, a surgeon does not seek a binary confirmation that a "suturing" or "stitching" event occurred, they require the precise localization of that event to evaluate procedural integrity or investigate post-operative complications. For instance, identifying the specific interval of a suture within 180 minutes of footage is a retrieval task that conventional models are unable to reliably accomplish.

Temporal grounding addresses this challenge by anchoring semantic understanding to a rigorous timeline, which is crucial for effectively addressing complex video understanding tasks \cite{univtg}.
This anchoring enables precise retrieval of specific moments by natural language query, causal linking of events separated by minutes, and structured summarization where every claim traces back to a verifiable segment. Without robust temporal grounding, the outputs of video-language models remain ungrounded, lacking the evidence required for applications where temporal accuracy is non-negotiable. By enforcing a strict mapping between semantic content and the temporal axis, we enable a deeper, more structured evidence-based form of video reasoning necessary for deploying VLMs in medical, legal, and industrial auditing. However,  without a benchmark that jointly evaluates answer correctness and temporal localization, there is no way to know whether any of these capabilities have actually been achieved.

% Temporal grounding addresses this challenge by anchoring semantic understanding to a rigorous timeline. \,This alignment facilitates three primary capabilities essential for professional-grade video models:
% \begin{enumerate}
%   \item High-Precision Retrieval: Enabling users to navigate massive datasets by pinpointing minute-long intervals based on complex natural language queries
%   \item Causal Reasoning: Establishing a structural backbone that allows the model to link separate events, such as a specific instrument choice at minute 14 and a subsequent physiological response at minute 47, thereby moving beyond correlation and towards a logical understanding of the workflow
%   \item Structured Summarization: Transforming hours of raw footage into a hierarchical, evidence-based report where every summarized claim is linked to a verifiable temporal segment
% \end{enumerate}

Current video benchmarks provide questions and answers but lack explicit temporal grounding. They do not direct models to the specific moments where answers can be found. This leaves models unable to reliably locate key events or reason precisely about when things occur in a video. While some works propose temporal understanding, they remain clip-level discrimination tasks that neither require models to localize events in time nor reason about their precise temporal position within a video, leaving temporal localization entirely unaddressed \cite{temporalbench}. Without the isolation of the temporal structure, models are left to rely on static visual cues or linguistic shortcuts. Meanwhile, traditional temporal grounding datasets \cite{charades} \cite{activitynetqa} use regression-based evaluation over short videos with known distributional biases, and are incompatible with the standardized multiple choice question (MCQ) evaluation paradigm that has become the lingua franca of VLM assessment. The problem is not that models answer incorrectly, but that we have no way to verify whether a correct answer reflects genuine temporal understanding or a lucky guess. The solution is conceptually simple: every correct answer should carry a verified timestamp, so that answer accuracy and temporal localization can be jointly evaluated.

% Some other works propose temporal ordering

% These benchmarks also often fail to capture the discrepancy between semantic recognition (understanding?) and temporal localization. More recent works suggest that VLM benchmark scores systematically overestimate video understanding performance on real-world tasks. TimeBlind \cite{timeblind} demonstrates that a state-of-the-art (SOTA) model like Gemini 3 Pro, despite scoring 88.4\% on Video-MME \cite{videomme}, achieves only 48.2\% instance accuracy on fine-grained temporal discrimination tasks in short 10-second clips, drastically trailing the human baseline of 98.2\%. This suggests that high performance on existing QA metrics does not guarantee a model's ability to anchor its reasoning in time. This discrepancy between QA accuracy and temporal precision is a byproduct of current evaluation construction. 
% Most existing benchmarks do not isolate temporal structure, leaving models to rely on static visual cues or linguistic shortcuts. 

To address this methodological gap, we introduce \textbf{ChronoQA}, a benchmark that directly measures a model's temporal grounding abilities. Each question is a timestamp-anchored MCQ benchmark drawn from long-form video (30–90 minutes). ChronoQA utilizes a dual-metric evaluation protocol that distinguishes models that answer correctly, from models that both answer correctly and identify accurate supporting evidence. 

Unlike existing benchmarks that rely on manual annotation or focus narrowly on perceptual recall, ChronoQA enforces strict temporal grounding without temporal leakage, ensuring questions must never reference time (such as "when," "at the beginning," "in the first segment"), yet every correct answer carries precise timestamp annotations locating the evidence within the video. By utilizing multimodal blind filtering and requiring timestamp-anchored answers, we eliminate the static shortcuts prevalent in existing suites and force the model to engage with the temporal axis of the video. This design ensures that model answer selections are evaluated on what they understand, not on their ability to pattern-match temporal cues in the question's text.

The benchmark is organized around a hierarchical taxonomy of five complementary reasoning categories: (i) Action Understanding, testing whether a model can identify actions, tasks, or behaviors occurring within a specific scene or setting; (ii) Needle-in-a-Haystack Retrieval, requiring precise recall of fine-grained visual and textual details (brand names, colors, clothing, on-screen text, counts, and spatial attributes) anchored to specific people, objects, or scenes distributed across the full video; (iii) Temporal Ordering, probing a model's understanding of the sequence and chronology of events, including step ordering, immediate predecessors/successors, and transition sequences between named states; (iv) Causal and Cross-Segment Reasoning, demanding integration of evidence across multiple disjoint video segments to infer causation, motivation, recurring patterns, and consequential relationships; and (v) Narrative Understanding, assessing holistic comprehension of the video's overarching narrative, primary subject, target audience, or overall tone, questions answerable only by a model that has processed the video in its entirety.

We evaluate several SOTA Multimodal Large Language Models (MLLMs), like Gemini 3 Pro, as well as strong open-source models like Qwen3-VL and Molmo2 on [Y] video-question pairs. We find....FINDINGS TBD

Our contributions are:

\begin{enumerate}
  \item ChronoQA, the first benchmark that requires models to both answer correctly and localize the supporting evidence in long-form video.
  % a [Y]-question timestamp-anchored MCQ benchmark that focuses on temporal grounding in long-form videos covering five temporal reasoning categories.
  \item A grounding integrity check that exposes models answering correctly for the wrong reasons, addressing a failure mode invisible to existing benchmarks.
  % A dual-metric evaluation protocol (QA Accuracy + mIoU) that differentiates between models that provide correct answers and those that not only deliver correct answers but also pinpoint relevant supporting evidence.
  \item An evaluation of several leading state-of-the-art (SOTA) MLLMs and their temporal grounding abilities.
\end{enumerate}


% VLMs have been shown capable of exploiting linguistic patterns. MMStar demonstrated
% that the initial Gemini Pro model \cite{gemini1.0} was capable of achieving 42.9\% on MMMU \cite{mmmu} without
% any visual context. The subsequent MMMU-Pro \cite{mmmu-pro} was created specifically to eliminate
% this problem of creating questions solvable by text, without visual context. In video
% benchmarks, models can frequently succeed using frame-level spatial cues without genuine
% temporal reasoning. MMBench-Video \cite{mmbenchvideo} found that some models (like Qwen-VL-Chat \cite{qwen-vl}), can perform
% better with a single frame of context rather than multiple, indicating that temporal
% information is not always necessary. ChronoQA's timestamp anchoring makes taking these
% kinds of shortcuts structurally infeasible, requiring grounding the correct answer to
% a specific moment in the video. We present three core contributions:



% Evaluation proceeds in two passes: answer selection (QA Accuracy) and temporal localization (Mean Temporal IoU), enabling diagnosis of whether a correct answer reflects genuine grounding or surface-level pattern matching. Our contributions are:


\todo{Add 2D scatter plot figure (Temporal Granularity vs Video Duration) positioning ChronoQA}
\todo{Add sample ChronoQA question figure showing video timeline, question text, five answer options each with paired timestamps, and the correct answer highlighted}


%%
%% ========================================================================
%%  RELATED WORK
%% ========================================================================
%%
% \section{Related Work}
% % Structured, critical review of VLM architectures, benchmarks, and temporal/video
% % understanding. Establishing where each falls short and how ChronoQA addresses that gap.

% \subsection{Vision-Language Models for Video}
% Current VLMs share a common architectural pattern: a vision encoder (SigLIP \cite{siglip}, InternViT \cite{internvl}, etc.) converts sampled frames into patch tokens, which are projected into the token space of an LLM backbome. This frame-sampling design is the primary bottleneck for temporal grounding, as inter-frame temporal signals are discarded, and fine-grained motion or transition information depends entirely on frame sampling density.

% Proprietary models attempt to address this constraint through larger context windows. Gemini 2.5 Pro supports approximately 6 hours of video sampled at 1 fps within a 2-million-token context length \cite{gemini2.5-video}, achieving 82\% on Video-MME. Gemini 3 Pro reduces this context window to 1 million tokens while improving multi-step reasoning. Unlike GPT-4, the natively multimodal GPT-5 scores 86.7\% on Video-MME, and up to 78.4\% on MMMU. Even with these expanded context windows, VLMs are shown to exhibit a degradation in performance as video duration increases \cite{longvideosurvey}, indicating that simply increasing context length is insufficient for resolving temporal grounding limitations. Open-source models like Qwen3-VL and Molmo2 introduce spatio-temporal grounding and video tracking, but operate at smaller effective temporal resolutions.

% While humans achieve over 98\% accuracy on SpookyBench \cite{spookybench}, fifteen state-of-the-art VLMs failed to identify any content, scoring 0\%. This categorical failure highlights a critical time-blindness in current architectures, which rely heavily on spatial features and struggle to decode information presented through purely temporal noise sequences, directly motivating ChronoQA's focus.

% Current state-of-the-art VLMs struggle with spatio-temporal awareness. VLM4D \cite{vlm4d} illustrates this, showing models can still fail at basic physical movement, often struggling to distinguish between clockwise/counter-clockwise rotations or simple lateral movements (e.g., a car moving right being labeled as moving left). This motivates a benchmark focused on temporal grounding, assessing how things change over time, rather than just spatial comprehension.

% While current state-of-the-art VLMs show promise in general video understanding, they fundamentally struggle with spatio-temporal awareness. As VLM4D \cite{vlm4d} illustrates, models frequently fail to resolve basic physical dynamics, such as distinguishing between clockwise and counter-clockwise rotations. This persistent time-blindness motivates the need for benchmarks like SpookyBench \cite{spookybench}, which isolates temporal grounding by stripping away all spatial shortcuts. By forcing models to derive meaning solely from changes over time, we can evaluate a model's inherent temporal reasoning independent of its spatial comprehension. While humans achieve over 98\% accuracy on SpookyBench, fifteen state-of-the-art VLMs failed to identify any content, scoring 0\%. This categorical failure further highlights the critical time-blindness in current architectures that struggle to decode information presented through purely temporal noise sequences, directly motivating ChronoQA's focus.
\section{Related Work}

\subsection{Video QA Benchmarks}
Current Video Question Answering (VideoQA) benchmarks largely overlook the necessity of fine-grained temporal grounding. For instance, Video-MME \cite{videomme} is leveraged extensively for evaluating frontier models, spanning 900 videos, 30 subjects and 2,700 questions. However, its temporal reasoning and event localization subtasks, among its 12 task categories, constitute a small fraction of the evaluation and none of its questions require the model to produce or verify timestamps or time intervals. SEED-Bench \cite{seedbench} covers action recognition, prediction, and procedure understanding across three video dimensions, but similarly lacks any timestamp precision requirement. EgoSchema \cite{egoschema} tests holistic comprehension over 3-minute egocentric clips using a 5-way MCQ format, and uses a temporal certificate metric to verify that correct answers require attending to the full clip rather than a single frame. However, holistic comprehension of a 3-minute clip is a different capability from timestamp-level grounding and precise localization in extended 30–90 minute durations. While this metric captures whether holistic attention to the clip is required, it does not require the model to produce or verify a timestamp prediction. ChronoQA extends this intuition to an explicitly grounded evaluation. Rather than certifying that a question requires the full clip to answer, every question in ChronoQA carries a verified timestamp interval that the model must identify alongside its answer selection.

A separate class of benchmarks evaluates temporal structures but does so at a coarse granularity. MVBench \cite{mvbench} treats action localization through broad beginning, middle, or end bins, effectively conflating temporal ordering with localization precision. NExT-QA \cite{nextqa} advances beyond descriptive benchmarks by explicitly targeting causal and temporal action reasoning across 5,440 videos, demonstrating that top-performing models achieve near-human accuracy on descriptive questions, but fail substantially on causal and temporal ones. However, NExT-QA videos average only 44 seconds in length, and no question requires the model to identify a supporting timestamp interval. Instead, the benchmark measures whether the model reasons correctly about temporal structure, not whether it can locate that structure within a longer video. Another work, TemporalBench \cite{temporalbench}, assesses sensitivity to temporal dynamics such as action frequency, motion magnitude, and event order across approximately 10K QA pairs, but does so via binary caption discrimination rather than timestamp anchoring. These benchmarks tackle temporal reasoning, understanding sequence and causality, while leaving the more rigorous problem of temporal localization largely unaddressed.

Among existing works, LongVideoBench \cite{lvbench} most closely aligns with the objectives of ChronoQA, featuring 6,678 MCQs for videos up to one hour. However, its "referring reasoning" format provides timestamps as input features, inverting the grounding task, prompting models to describe events at a given interval rather than identifying the interval corresponding to a described event.

% Finally, traditional grounding datasets such as Charades-STA \cite{charades} and ActivityNet Captions \cite{activitynetqa} typically utilize regression-based start/end timestamp prediction. This methodology is structurally incompatible with the MCQ paradigm and is prone to distributional biases. Notably, Charades-STA exhibits a high density of ground-truth moments near video boundaries, while ActivityNet moments frequently span over 30\% of the video duration, rewarding imprecise predictions \cite{tgsv}. ChronoQA departs from these regression-based frameworks, utilizing a contrastive MCQ structure that necessitates the discrimination of temporally plausible alternatives.

\begin{table}[h]
  \centering
  \caption{Comparison of ChronoQA with existing benchmarks.}
  \label{tab:benchmark_comparison}
  % \todo{Add comparison table}
  % Columns: Benchmark, Venue, Modality, \# Videos, \# Questions, Video Duration,
  %          Answer Format, Timestamp-Anchored?, Temporal Grounding Focus?, Scale Pipeline
  % Rows: Video-MME, MVBench, EgoSchema, TemporalBench, LongVideoBench,
  %        SEED-Bench, Charades-STA, ActivityNet Captions, ChronoQA
\end{table}

\subsection{Temporal Reasoning in Vision-Language Models}
Current VLMs share a common architectural bottleneck: a vision encoder converts sampled frames into patch tokens projected into an LLM backbone, discarding inter-frame temporal signals entirely. Fine-grained motion and transition information therefore depends on frame sampling density, which remains sparse even in frontier models. Proprietary systems have attempted to compensate through larger context windows. Gemini 2.5 Pro supports approximately six hours of video within a two-million-token context, and the natively multimodal GPT-5 achieves 86.7\% on Video-MME, yet even with these expanded budgets, VLMs exhibit systematic performance degradation as video duration increases, indicating that scaling context length alone does not resolve the underlying temporal grounding deficit \cite{longvideosurvey}. Open-source models like Qwen3-VL and Molmo2 introduce spatio-temporal grounding and video tracking, but operate at smaller effective temporal resolutions.

This architectural limitation has measurable consequences for model behavior. TempCompass \cite{tempcompass} evaluates temporal perception across five aspects: action, speed, direction, attribute change, and event order. Using conflicting video pairs that share identical static content but differ along a single temporal dimension, they neutralize single-frame shortcuts and language priors. Across eight SOTA Video LLMs, models perform near-randomly on speed, direction, and attribute change, achieving reasonable accuracy only on action, a category largely solvable from static features. NExT-QA \cite{nextqa} demonstrates a parallel trend in reasoning, where the disparity between model and human performance is significantly greater for causal and temporal questions than for descriptive ones, even when dealing with brief and well-structured videos. Categorical failures appear even on simpler temporal tasks. SpookyBench \cite{spookybench} presents information exclusively through temporal noise sequences, stripping away all spatial content, resulting in fifteen SOTA VLMs scoring 0\% on tasks where humans achieve over 98\% accuracy. Another work, VLM4D \cite{vlm4d}, shows that models frequently fail to resolve basic physical dynamics, such as distinguishing clockwise from counter-clockwise rotation or correctly labeling the direction of lateral movement. Together, these results suggest that current VLMs exploit static visual shortcuts wherever available and fall back towards chance when those shortcuts are removed.
% Taken together, these results indicate that current models have not learned to derive meaning from temporal change at any level of abstraction — from low-level motion direction to high-level causal reasoning.

Critically, both TempCompass and NExT-QA operate on short clips where temporal evidence is local and the search window is small. Whether this failure pattern holds in long-form video remains unanswered. Existing temporal grounding datasets such as Charades-STA\cite{charades} and ActivityNet-QA\cite{activitynetqa} use regression-based evaluation but are similarly confined to short videos with known positional biases. Only 40\% of ActivityNet moments span more than 30\% of the video duration, allowing models to exploit distributional shortcuts rather than performing genuine localization. No existing benchmark jointly measures answer accuracy and timestamp precision in a long-form setting. ChronoQA addresses this gap with a dual-metric evaluation protocol that reports both QA accuracy and mIoU over predicted intervals on videos of 30 to 90 minutes, making explicit that aggregate answer accuracy without grounding verification cannot be taken as evidence of temporal grounding ability.





%%
%% ========================================================================
%%  DATA COLLECTION & ANNOTATION
%% ========================================================================
%%
\section{Data Collection \& Annotation}
% Describe the fully automated pipeline; sourcing videos, QA pair generation, quality filters,
% with sufficient detail for reproducibility.

\subsection{Video Sourcing \& Selection Criteria}

YouTube videos provide a large-scale, real-world source of long-form content with diverse metadata and content types. We perform large-scale and filtering by duration (30-90 minutes), content type (instructional, narrative, procedural, sports, vlogs), licensing (Creative Commons or fair-use) and metadata availability. A visual diversity filter is applied to ensure  sufficient temporal dynamics for meaningful temporal reasoning questions. We compute a "delta" score for each video, representing the average visual change across sampled frames using CLIP embeddings and cosine similarity. Videos with low delta scores (indicating static content) are excluded.
% Contrast with Video-MME's \cite{videomme} manual curation of 900 videos, ChronoQA achieves \textbf{44$\times$ the video count} through automation while maintaining quality via multi-stage filtering.

\subsection{Caption Generation Pipeline}

In order to obtain precise event timestamps, we segment each video into fixed-length clips of $10$ seconds. Each segment is passed to Qwen3VL, which samples frames at $1$ frame per second, and projects them into a pixel-budget space of $T_{Video} = C \cdot 32 \cdot 32$ tokens, where $C$ is a user-specified token budget per sampled frame.

% TODO: Section about the importance of segment lengths? We looked at 10s, 20s, and 30s.
% 5s probably too expensive, 20s and 30s had their own caveats. Refer to slides for more info.

Ten-second clips provide sufficient temporal resolution to capture fine-grained event details without over-generalizing across broader scene transitions.

% We segment videos into 10-second clips, which provides sufficient temporal resolution to capture fine-grained event details without over-generalizing across broader scene transitions. 

% We experimented with different fixed-length clips (10s, 20s, and 30s) to understand the impact of temporal granularity on caption quality and downstream question generation. The choice of segment length affects the level of detail in the captions, with shorter segments (like 10s) providing more granular descriptions of specific moments, while longer segments (like 30s) capture broader events and themes. We found that 10-second segments strike a good balance, providing sufficient temporal resolution to capture distinct events without overwhelming the model with too much information in a single segment. Longer segments tended to produce more generalized captions that missed finer details, while shorter segments increased computational costs without significant gains in caption quality.


\subsubsection{Caption Token Budget}

Qwen3-VL processes images with a SigLip-2 Vision Transformer \cite{siglip2}, converting pixels into a dynamic number of tokens before compressing them spatially with a learned two-layer MLP merger, and then injecting into the language model at multiple depths using Deepstack \cite{deepstack}. The visual token count is entirely determined by the input image resolution. An image of $H \times W$ pixels will produce exactly $\left(\frac{H}{32} \times \frac{W}{32}\right)$ visual tokens after processing. The factor of $32$ arises from two stages of spatial compression. Firstly, the $16 \times 16$ patch embedding, which is followed by $2 \times 2$ spatial merging from the MLP merger. A $640 \times 640$ image yields roughly $20 \times 20 = 400$ visual tokens, a $1280 \times 1280$ image yields $40 \times 40 = 1{,}600$ tokens.

This budget is controlled during preprocessing stage by specifying \texttt{min\_pixels} and \texttt{max\_pixels}. The processor resizes images to maintain aspect ratio while keeping total pixel count within this range. For video, setting \texttt{total\_pixels} to cap the combined token count of all frames, and a temporal patch size of 2 provides a $2x$ temporal compression by merging consecutive frame pairs. Each video frame is divided into non-overlapping $32 \times 32$ pixel patches. 

The total visual pixel budget is:
\begin{equation}
  P_{video} = T_{Video} \times 32^{2} = T_{Video} \times 1024
\end{equation}

This cap is distributed across all sampled frames. For a $t$-second segment sampled at $2$fps, $F=2t$ frames yield $F / 2=t$ frame pairs after temporal compression, with each frame receiving $T_{FP} = T_{Video}/t$ tokens. Increasing $T_{Video}$ yields finer visual detail but raises two costs. Firstly, the token sequence sent to the model grows proportionally, increasing both memory pressure and inference time. Secondly, a larger token budget per clip expands the KV cache allocated per sequence, directly constraining concurrent batch throughput before PagedAttention \cite{pagedattention} must begin preempting sequences. If $T_{Video}$ is too low relative to $t$, each frame receives so few tokens that fine visual details are lost.

All segments inputs are submitted concurrently to the vLLM inference engine. Qwen3-VL-30B generates one free-form caption per segment, which is recorded alongside its timestamp and source video ID.

\subsection{Methods}

\subsubsection{Overview}

We present an automated pipeline for generating temporally-grounded multiple-choice questions from equal duration video narration captions. The pipeline operates in four stages:
\begin{enumerate}
  \item Question-Answer Generation
  \item Validation and Cleaning
  \item Distractor Synthesis
  \item Blind Filtering
\end{enumerate}

All LLM inference was done using the vLLM inference engine with Pydantic-enforced structured outputs to guarantee well-formed JSON while minimizing post-hoc parsing.

\subsubsection{Stage 1: Question-Answer Generation}

The pipeline takes as input a JSON file mapping video IDs to caption data, where each entry contains a list of timestamped sentences and the original video's duration. We define five specialized prompt templates, each targeting a distinct facet of video understanding:

\begin{itemize}
  \item \textbf{Action Understanding}: Tests comprehension of actions or tasks within specific scenes, anchored by scene/setting descriptions (e.g., "What is the character next to the chair doing in the outdoor setting?").
  \item \textbf{Needle-in-a-Haystack}: Retrieves specific visual or textual details such as colors, brands, clothing, or on-screen text, requiring contextual anchors to uniquely identify subjects.
  \item \textbf{Ordering}: Tests knowledge of action sequences using relative ordering language (like "immediately after" or "right before") rather than absolute positional references.
  \item \textbf{Causal Understanding}: Requires cross-segment reasoning about causation and motivation, connecting evidence from two or more separate parts of the video.
  \item \textbf{Narrative Understanding}: Tests holistic comprehension of the overarching narrative, primary subject, core message, or genre.
\end{itemize}

Each template generates 3 questions per video. The model is instructed to follow the prompt guidelines and use the captions and their timestamps to generate a question-answer pair grounded in the video content. All templates are constrained to make questions that never reference time directly (no "when," "beginning," or timestamp references), and only the correct answer may carry a temporal annotation in the form of a start and end timestamp for the relevant segment. This ensures that questions can be presented to a viewer without having to reveal any temporal cues.

We use Qwen3-Next-80B-A3B-Thinking \cite{qwen3} constrained to a Pydantic schema via vLLM's \texttt{StructuredOutputsParams} to guarantee each response contains well-formed question-answer pairs with associated timestamps. However, QA generation can still fail for a variety of reasons. In some cases, the model hallucinates, exceeding the provided output token budget. In other cases, the model may return JSON with incorrect fields or incorrect formatting. Accounting for these cases, QA generation succeeds approximately 98\% of the time.

\subsubsection{Stage 2: Validation and Cleaning}

Generated QA pairs then undergo a series of quality control steps.

Firstly, we use pattern matching to detect questions that inherently require temporal information to answer (e.g., "When was X shown?") and discard these questions, as they would be unanswerable without timestamps in the answer. We prompt against the generation of questions like these, but apply this step to ensure that no such questions make it through.

Secondly, we apply timestamp merging to questions with multiple entries in the timestamp fields. For these questions, adjacent timestamp segments that are separated by less than 2 seconds are merged into a single timestamp range to avoid fragmentation. This was implemented because generated QA pairs would sometimes have multiple timestamp segments that are better represented as a single continuous range.

Finally, we expand the recorded timestamp range to be larger than the ground truth, to increase timestamp variability. Due to the $N$ clips of equal length that video captions are created with, the LLM will usually set ground truth timestamps in $10$ second chunks. As a result, generated ground truth timestamps are usually only as fine and as coarse as the $10$ second clip length. To create questions with more varied timestamp ground truths, we apply a ``timestamp jitter.'' This jitter records the starting timestamp a $\text{random}(1, 3)$ seconds earlier, and the ending timestamp a $\text{random}(1, 3)$ seconds later, than what is recorded in the timestamp ground truth returned by the LLM. This also provides tolerance for imprecise annotations returned by the LLM when timestamp ground truth ranges are smaller than the $10$ second segment length.

\subsubsection{Stage 3: Distractor Generation}

For each question, we generate four plausible but incorrect answer options (distractors) in a separate LLM pass. Within the prompt, we instruct the model to produce options that (a) are plausible within the question context, (b) belong to the same semantic category as the correct answer, (c) match the correct answer in length and casing, and (d) are unambiguously incorrect.

When relevant caption data is available, we extract narration segments within a $\pm 30$-second margin around the correct answer's timestamp and provide this localized context to the model, in order to ground distractors in actual video content. Questions for which fewer than 4 valid distractors are successfully generated are discarded. The correct answer and four distractors are then randomly shuffled, and the position of the correct answer is recorded as a letter (A-E).

\subsubsection{Stage 4: Blind Filtering}

Motivated by TemporalBench \cite{temporalbench}, we observe that language models can sometimes identify the correct answer without visual input by exploiting shared linguistic patterns between the question and the ground truth, particularly when distractors are lightly edited variants of the correct answer. To filter out questions susceptible to this bias, we apply a blind filtering procedure in which each QA sample, consisting only of the question text and five answer options, with no video, is presented to [X] state-of-the-art LLMs: GPT-5.4  and Qwen3-Next-80B-A3B-Thinking. Each model is prompted to select the correct answer under the same task framing used in full evaluation. Since the random baseline for a five-way multiple-choice question is 20\%, any sample where at least one model selects the correct answer is considered language-solvable and is removed from the dataset, ensuring that retained questions cannot be resolved through linguistic pattern-matching alone and require genuine visual grounding.

\subsubsection{Scale Statistics}

Each final QA instance contains a video ID, question text, correct answer letter, and five answer options (each with associated timestamps), formatted as structured JSON suitable for use in video understanding benchmarks.

\todo{Final dataset: 40,000+ source videos to TBD evaluation subset (likely 5,000-10,000 curated questions) after filtering. Report statistics on video duration distribution, question-type distribution, temporal precision distribution, and domain diversity.}

\todo{Dataset Statistics table (Total Videos, Total Questions, Mean/Median/Max Video Duration,
       Questions per Task Category, Domain Distribution, Mean Temporal Precision of Answers)}
\todo{Distribution Plots (Video Duration, Answer Timestamp Distribution, Distractor Temporal Distance)}
\todo{Pipeline Diagram flowchart (YouTube Crawl -> Metadata Extraction -> Temporal Event
       Identification -> LLM-based QA Generation -> Quality Filtering Cascade -> Human
       Validation -> Final ChronoQA Dataset). Annotate each stage with approximate counts.}


       %%
%% ========================================================================
%%  BENCHMARK DESIGN & TAXONOMY
%% ========================================================================
%%
\section{Benchmark Design}
% Present ChronoQA task structure, question types, and design principles, showing how each
% design decision is motivated by the limitations of prior work.

\subsection{Task Formulation}

Given a long-form video $V$ and a natural language question $Q$ that references a temporal event, transition, or property of the video, the model must select the correct answer $a_{i}$ from a set of \textbf{five} options $\{a_1, \ldots, a_5\}$. Each option consists of a \textbf{textual description} indicating the segment of the video that description refers to.

This task requires the model to:
\begin{enumerate}
  \item Interpret the contextual reference in $Q$ to identify the relevant moment(s) of the video.
  \item Inspect the corresponding video segment(s).
  \item Select the answer whose textual description correctly matches the visual content of that segment.
\end{enumerate}

Incorrect options are wrong because their descriptions are visually false \textbf{at that particular moment in time}. This design ensures that contrastive reasoning of answer options must be done to eliminate distractors, requiring the model to perform genuine visual grounding. A model without visual context seeing five plausible descriptions based on information present in the video has no basis for discrimination beyond linguistic guessing, achieving near-random performance (approximately 20\%) on well-constructed questions.

\subsection{Task Taxonomy}

\textbf{ChronoQA} organizes its questions into five categories, each targeting a distinct capability. These categories are roughly ordered by the scope of video context required, from segment-based visual retrieval to full-video narrative comprehension.

\subsubsection{Action Understanding} \newline

\textbf{Goal:} Identify what action is occurring at a specific moment in the video, including procedural details on how the tasks or activities are being carried out.

Action Understanding questions assess a model's ability to accurately characterize a video segment's content, understanding how it unfolds, not just detecting that an action is occurring. This requires fine-grained visual perception of things like tool usage, body posture, object manipulation, and sequential steps within a single continuous action. Questions in this category are anchored to a specific segment and do not require cross-segment reasoning.

\textbf{What it addresses:} Existing benchmarks like MVBench \cite{mvbench} test coarse action recognition (e.g., "is the person running or walking?") but rarely probe procedural details. \textbf{ChronoQA}'s Action Understanding questions demand the level of detail a domain expert would notice. For example, identifying the difference between two similar techniques performed at different moments in a tutorial video.

\todo{Add example questions for Action Understanding.}

\subsubsection{Needle-in-the-Haystack} \newline

\textbf{Goal:} Retrieve a specific visual or textual detail like an object's characteristics, brand names, colors, materials, on-screen text, or object counts, from anywhere within the full video timeline.

Needle-in-the-Haystack questions test a model's ability to locate and report precise details that only appear briefly on-screen. Unlike the other categories, these questions are explicitly constructed so that the correct-answer timestamps are uniformly distributed across the video's full duration. This means that the detail concerned by the question is equally likely to appear either early, in the middle, or near the end. This design choice addresses the well-documented start/end bias of previous temporal grounding datasets such as Charades \cite{charades}, where segments are clustered disproportionately near the video's boundaries, enabling models to take advantage of positional priors instead of performing genuine information retrieval.

The model must be capable of sustaining attention across the entire video to locate the target detail, and cannot rely on narrative context to narrow the search window.

\textbf{What it addresses:} Most video benchmarks either test contextually prominent events (that are therefore easy to anticipate) or sample questions from predictable temporal positions. Needle-in-the-Haystack isolates retrieval capabilities from reasoning capabilities, creating a pure test of a model's ability to find a specific detail in a long, unstructured visual stream.

\todo{Add example questions for Needle-in-the-Haystack.}

\subsubsection{Ordering} \newline

\textbf{Goal:} Reason about the sequential structuring/ordering of events and actions in the video. What happens before or after a specific event, how transitions occur from scene to scene, and in what order items or steps are used and/or encountered.

Ordering questions test the ability to temporally sequence events. The model must understand what occurs at a given moment, as well as how that moment relates to adjacent events in time. This includes reasoning about direct relationships like before/after, identifying transitions between scenes or activities, and determining the order in which objects are used or steps are taken within a process.

These questions differ from Action Understanding because they are inherently relational. The correctness depends on the position of an event relative to other events, not just the local content. This differs from Causal Understanding because it does not require inferring why those events occur in a particular order, only that they do occur in that order and when they do.

\textbf{What it addresses:} MVBench \cite{mvbench} includes sequencing-based tasks but primarily uses short clips where temporal order can often be inferred from a handful of frames. \textbf{ChronoQA}'s ordering questions span long-form videos where events may be separated by spans of several minutes, requiring the model to maintain a coherent temporal mapping of the video rather than simply relying on local context.

\todo{Add example questions for Ordering.}

\subsubsection{Causal Understanding} \newline

\textbf{Goal:} Integrating evidence from two or more non-adjacent segments of the video to understand the motivations of characters, cause-and-effect relationships, and patterns that only become apparent when examined across different scenes.

Causal Understanding is the most demanding reasoning category of \textbf{ChronoQA}. Every question in this category is constructed such that \textbf{no single video segment contains sufficient information to answer correctly}, and the model must synthesize observations from multiple, temporally separate parts of the video and draw an inference that connects them. This may involve understanding a character's motivations for doing particular actions, identifying what earlier event caused a later outcome (causality), or recognizing a recurring pattern whose significance only emerges through repetition across scenes (pattern recognition).

This category establishes a distinction between \textbf{comprehension} (understanding what happens at a particular moment in time) and \textbf{reasoning} (understanding why it happens, given what came before and after). It is also the most resistant to frame-sampling shortcuts. A model that processes only a subset of frames may observe the effect without ever seeing the cause, or vice versa, but will struggle to make connections between the cause and effect without seeing the other.

\textbf{What it addresses:} EgoSchema \cite{egoschema} requires holistic reasoning across 3-minute clips, but its questions are answerable from the general arc of the clip. LongVideoBench \cite{lvbench} tests referring reasoning, but supplies relevant timestamps in the question. Neither requires the model to independently identify and integrate evidence from multiple non-adjacent segments. This category makes cross-segment synthesis a structural requirement.

\todo{Add example questions for Causal Understanding.}

\subsubsection{Narrative} \newline

\textbf{Goal:} Assess high-level synthesis of a video's overarching narrative arc, primary subject matter, core message, or overall mood and tone. This requires the model to have processed the video in its entirety.

Narrative questions test whether a model has formed a coherent global understanding of the video, and not just a collection of local observations. Correct answers require understanding a video's broader context, including its genre, overarching goal, main thesis, primary technique, or dominant tone. Questions are anchored to properties that can only be determined from the full video. A question that can be answered with the first 30 seconds of the video as context, or from any single representative segment, would fail to qualify.

Despite testing high-level understanding, Narrative questions are held to strict objectivity standards. Questions that are too vague to score reliably (like \emph{"What is the mood of the video?"} without a verifiable/objective correct answer) are excluded during quality filtering. All Narrative questions must have a single defensible correct answer that a human viewer would agree upon after watching the complete video.

\textbf{What it addresses:} SEED-Bench includes "action prediction" and "procedure understanding" as high-level video tasks, but these tasks focus on fine-grained categories of basic actions and action segmentation. Video-MME \cite{videomme} includes "summarization" as a task category, but answers are often derivable from a few key scenes. \textbf{ChronoQA}'s Narrative category specifically enforces full-video dependency. The correct answer should reflect the \emph{aggregate} of the video's content, not any one segment of the video.

\todo{Add example questions for Narrative.}

\subsection{Design Principles}

\begin{itemize}
  \item \textbf{Temporal grounding prevents linguistic shortcuts} (addresses MMStar's \cite{mmstar} finding that 42.9\% of MMMU is solvable text-only): All answer options reference the same video segment, so models cannot eliminate distractors by comparing timestamps. Only genuine visual inspection of the segment can distinguish correct from incorrect descriptions.
  \item \textbf{Long-form videos prevent single-frame solving} (addresses MMBench-Video's \cite{mmbenchvideo} finding that 1-frame sometimes matches 8-frame performance): Questions require reasoning over events separated by minutes, not seconds, preventing single-frame solving.
  \item \textbf{Balanced answer distributions prevent position bias} (adopts MMBench's \cite{mmbench} CircularEval insight): Correct answers are uniformly distributed across positions $A, B, C, D, E$.
  \item \textbf{Five-option MCQ enables standardized evaluation} (addresses Charades-STA/ActivityNet Captions' reliance on IoU-based regression metrics): ChronoQA uses accuracy, the universal metric of VLM benchmarks, while encoding temporal precision through timestamp pairing.
  \item \textbf{Distractor design philosophy}: Distractors are not random. Instead, they are semantically plausible events that occur at different timestamps in the same video, or the correct event paired with an incorrect timestamp. This ensures the task genuinely requires temporal grounding, not just event recognition.
\end{itemize}

\todo{Task Taxonomy Visualization figure (hierarchical tree showing 5 task categories with examples)}
\todo{Distribution Visualization figure (questions across task categories and video duration ranges)}
\todo{Distractor design illustration figure}


%%
%% ========================================================================
%%  EVALUATION PROTOCOL
%% ========================================================================
%%
\section{Evaluation Protocol}

\subsection{Evaluated Models}

ChronoQA evaluates seven models spanning the full range of capabilities in the current VLM landscape, including both closed-source and open-source models.

\begin{table*}
  \caption{Models evaluated on ChronoQA.}
  \label{tab:models}
  \begin{tabular}{llll}
    \toprule
    Model & Provider & Access & Context / Frame Limit \\
    \midrule
    GPT-4o & OpenAI & API & $\sim$128K tokens \\
    GPT-5 & OpenAI & API & 400K tokens \\
    Gemini 1.5 Pro / 2.0 / 3.0 & Google & API & Native long-video; up to 1M tokens \\
    Claude 3.5 Sonnet / 3.7 & Anthropic & API & $\sim$200K tokens; image-based video input \\
    LLaVA-NeXT & Open-source & Local inference & ~4k tokens \\
    Qwen3-VL &  Open-source & Dynamic resolution & 256K (expandable to 1M) \\
    Molmo2 & Allen AI / Open-source & Local inference & ~36K \\
    \bottomrule
  \end{tabular}
\end{table*}

All proprietary models are evaluated through their respective public APIs at the most capable available version at the time of paper writing. Model versions and API snapshot dates are recorded and reported for reproducibility. Open-source models are evaluated using their official HuggingFace checkpoints at their default recommended inference configuration.

\subsection{Input Specification}

We provide each model with a standardized input consisting of three components:
\begin{enumerate}
  \item \textbf{Video Input:} The source video, formatted as uniformly sampled frames at a rate standardized to a practical input limit for the model. For models with native long-video ingestion (e.g., Gemini), the full video is passed directly. For all other models, uniform frame sampling is applied.
  \item \textbf{Question Text:} The natural language question $Q$.
  \item \textbf{Answer Options:} Five labeled options (A-E), each consisting of a textual description and its associated timestamp interval $[t_{start}, t_{end}]$.
\end{enumerate}

The full prompt template is fixed across all models. No chain-of-thought instructions are included in the primary evaluation (zero-shot direct answer).

\subsubsection{Frame Sampling Budget}

To ensure fair comparison between models with different input limits, we defined a standardized frame budget of 64 uniformly sampled frames as the default evaluation setting. Gemini models are evaluated at both the $64$-frame budget and with native full-video ingestion, to isolate the effects of frame density from architectural long-video support.

\subsection{Metrics}

ChronoQA employs two complementary evaluation metrics targeting distinct aspects of a model's capabilities. One measures the accuracy of the final answer selection, and the other measures the quality of temporal grounding that underlies correct answers. This dual-metric design is directly motivated by \cite{cgbench}, which demonstrated that MCQ accuracy alone is insufficient to verify that correct answers are genuinely grounded. A model may select the right answer through surface-level reasoning without ever identifying the relevant video segment.

\subsubsection{QA Accuracy (Primary Metric)}
QA Accuracy is the percentage of questions for which the model selects the correct answer option:
\begin{equation}
  \text{Acc} = \frac{1}{N} \sum_{i=1}^{N} \mathbf{1}[\hat{a}_i = a_i^*]
\end{equation}

where $N$ is the total number of questions, $\hat{a}_i$ is the model's selected option, and $a_i^*$ is the ground-truth correct option.

QA Accuracy is reported at two levels of granularity:
\begin{itemize}
  \item \textbf{Overall Accuracy}: Computed across all questions in the dataset.
  \item \textbf{Per-Category Accuracy}: separately for each of the five task categories (Action Understanding, Needle-in-the-Haystack, Ordering, Causal Understanding, Narrative).
\end{itemize}


The random baseline for a 5-option MCQ is 20\%. A text-only baseline (no video input) must achieve below 20\% to confirm that linguistic patterns in the question and answer text do not disambiguate the correct option.

\subsubsection{Mean Temporal IoU (mIoU) - Grounding Metric}
mIoU measures the quality of a model's temporal localization. Given a question, how accurately does the model identify the relevant video segment, regardless of whether it selects the correct answer?
Following the standard temporal grounding evaluation protocol, temporal IoU (tIoU) for a single prediction is defined as:

\resizebox{\linewidth}{!}{
  \begin{equation}
    $$
    \text{tIoU} = \frac{
      \sum_{i \in \mathcal{G},\, j \in \mathcal{P}} \max\bigl(0,\, \min(b_i, d_j) - \max(a_i, c_j)\bigr)
    }{
      \sum_{i \in \mathcal{G}} (b_i - a_i)
      + \sum_{j \in \mathcal{P}} (d_j - c_j)
      - \sum_{i \in \mathcal{G},\, j \in \mathcal{P}} \max\bigl(0,\, \min(b_i, d_j) - \max(a_i, c_j)\bigr)
    } \times 100\%
    $$
  \end{equation}
}

where $a_{i}, b_{i}$ are the start and end timestamps of the $i$-th ground truth interval of $G$. $c_{j}, d_{j}$ are the start and end timestamps of the $j$-th predicted interval of $P$. We calculate the mean IoU (mIoU) by averaging the tIoU scores obtained by the model across all question queries. 
\begin{equation}
  \text{mIoU} = \frac{1}{N} \sum_{i=1}^{N} \text{tIoU}_{i}
\end{equation}

\paragraph{\textbf{Eliciting Timestamp Predictions.}}
Since ChronoQA is an MCQ benchmark and most evaluated models are not natively trained for timestamp regression, timestamp predictions are obtained with a structured prompting strategy. After selecting an answer option, models are further prompted to output the specific timestamp interval within the video that they believe best supports their answer (e.g., "State the start and end time in seconds of the video segment that directly answers this question."). This protocol follows the white-box evaluation paradigm of CG-Bench, which requires models to explicitly output grounding intervals alongside their MCQ selections.

For models that refuse, or fail, to produce a parsable timestamp (e.g., outputting a range covering the full video or a null response), a fallback interval of $[0,T]$ (full video duration) is assigned, yielding a tIoU score proportional to the fraction of the video covered by the ground-truth interval. This penalizes uninformative outputs without zeroing them entirely.

\paragraph{\textbf{mIoU as a grounding integrity check.}}
mIoU serves a diagnostic function beyond ranking, allowing us to detect cases where a model selects the right option but provides a grounding interval inconsistent with the correct timestamp, or a "lucky correct" answer. A model with high QA Accuracy but low mIoU is selecting correct answers for the wrong reasons, which is precisely the failure pattern CG-Bench was designed to surface. We additionally report Grounded Accuracy ($G-Acc$), which is the percentage of questions where the model both selects the correct answer and achieves $tIoU \geq 0.5$ on the grounding interval, providing a joint measure of answer correctness and temporal grounding quality.

\paragraph{\textbf{Recall@IoU (Supporting Metric).}}
Following CG-Bench, in order to improve the robustness of question grounding evaluation, we additionally report Recall@IoU, measuring the probability of successfully recalling clue intervals at various IoU thresholds. We calculate the average recall rate at IoU thresholds of $0.1, 0.2, 0.3, 0.4,$ and $0.5$ to determine the final result. This metric captures the distribution of grounding quality across the full IoU spectrum, providing a richer picture than mIoU alone, particularly for distinguishing models that frequently achieve partial overlap ($IoU ~0.3$) from those that more often achieve precise localization ($IoU ~0.5+$).

\subsection{Answer Extraction}
Reliable answer extraction is essential for consistent scoring across models with varying output styles. We adopt a two-stage extraction procedure, following MMBench:

Rule-based extraction (primary): Parse the model's output for an explicit option label ($A, B, C, D,$ or $E$) using regex matching. This succeeds for the majority of well-behaved model outputs.
LLM-based extraction (fallback): For outputs that do not contain an unambiguous option label, a separate GPT-4o judge is prompted with the model's full response and the original answer options to identify the intended selection. This handles cases where a model re-states the answer text rather than the option letter.

Timestamp extraction follows a similar two-stage approach, starting with regex parsing for structured outputs (e.g., "Start: 34.5s, End: 42.1s"), then followed by an LLM-based parser for free-form outputs.
If a model explicitly refuses to answer, or if both extraction stages fail, the question is scored as incorrect (0 accuracy) and assigned a full-video fallback interval for tIoU computation. Refusal rates are reported per model as a diagnostic.

\subsection{Baseline Configurations}
Four baseline conditions are evaluated to establish lower bounds and verify benchmark integrity. All baselines use the same answer extraction pipeline.

\begin{table*}
    \centering
    \begin{tabular}{llll}
        \toprule
        \textbf{Baseline} & \textbf{Input} & \textbf{Purpose} & \textbf{Target} \\ 
        \midrule
        Random & None & Theoretical lower bound & 20\% \\ 
        \addlinespace
        Text-only & Question + options (no video) & Confirm linguistic shortcut resistance & $<$ 20\% \\ 
        \addlinespace
        Single-frame & 1 random frame + question & Confirm temporal multi-frame requirement & $<$ 30\% \\ 
        \addlinespace
        Human & Full video + question & Establish performance ceiling & $\sim$85\% \\ 
        \addlinespace
        Audio/ASR-only & Transcript + question (no frames) & Determine auditory impact & Report as-is \\
        \bottomrule
    \end{tabular}
    \caption{Baseline Evaluation Metrics}
    \label{tab:baselines}
\end{table*}


The text-only baseline is the most critical integrity check: if any text-only model exceeds 20\% accuracy overall, the affected questions are flagged for manual review and potential removal from the evaluation set, following MMStar's multi-modal gain filtering procedure.

\subsection{Reproducibility and Access}
All evaluation code, prompt templates, answer extraction scripts, and model inference configurations are released publicly.

%%
%% ========================================================================
%%  EXPERIMENTS & RESULTS
%% ========================================================================
%%
\section{Experiments \& Results}
% Present comprehensive results across proprietary and open-source VLMs to demonstrate the
% benchmark's discriminative power and reveal systematic weaknesses in temporal grounding.

\todo{Complete section with the following subsections.}

\subsection{Main Results}
\todo{Main results table across all evaluated models.}

\subsection{Baseline Integrity Results}
\todo{Results from blind filtering showing shortcut resistance.}

\subsection{Performance by Video Duration}
\todo{Analysis of model accuracy as a function of video length.}

\subsection{Performance by Task Category}
\todo{Breakdown of results across the five task categories.}

\subsection{Performance by Temporal Precision}
\todo{Analysis of accuracy vs.\ temporal precision of answers.}

\subsection{Cross-benchmark Correlation}
\todo{Correlation of ChronoQA scores with existing benchmark scores.}


%%
%% ========================================================================
%%  ANALYSIS & INSIGHTS
%% ========================================================================
%%
\section{Analysis \& Insights}

\todo{Complete this section with:}

\subsection{Error Taxonomy and Qualitative Analysis}
\todo{Categorize common failure patterns with qualitative examples.}

\subsection{Frame Sampling Sensitivity Analysis}
\todo{Analysis of how frame sampling rate affects model performance.}


%%
%% ========================================================================
%%  CONCLUSION
%% ========================================================================
%%
\section{Conclusion}

\todo{Write conclusion summarizing contributions, key findings, limitations, and future work.}


%%
%% The acknowledgments section is defined using the "acks" environment
%% (and NOT an unnumbered section). This ensures the proper
%% identification of the section in the article metadata, and the
%% consistent spelling of the heading.
\begin{acks}
\todo{Add acknowledgments.}
\end{acks}


%%
%% Print the bibliography
%%
\printbibliography[title={References}]

\end{document}
